% Please do not change %
\documentclass[11pt]{article} % font size
\usepackage[margin=1in]{geometry} % margins
\usepackage{amsmath,amsthm,amssymb} % for the  common "math symbols"
\usepackage{algorithm}
\usepackage{algpseudocode}
\usepackage[shortlabels]{enumitem} % for enumeration
\usepackage{booktabs} % if you need tables
\usepackage{listings}
\usepackage{color}
\usepackage{quiver}
\DeclareMathOperator*{\argmax}{argmax}
\DeclareMathOperator*{\argmin}{argmin}

% \theoremstyle{definition}
\newtheorem{theorem}{Theorem}[section]
\newtheorem{corollary}[theorem]{Corollary}
\newtheorem{lemma}[theorem]{Lemma}
\newtheorem{proposition}[theorem]{Proposition}
\newtheorem{conjecture}[theorem]{Conjecture}
\theoremstyle{definition}
\newtheorem{definition}[theorem]{Definition}
\theoremstyle{definition}
\newtheorem{fact}[theorem]{Fact}
\theoremstyle{definition}
\newtheorem{example}[theorem]{Example}

\lstset{
  frame=none,
  xleftmargin=2pt,
  stepnumber=1,
  numbers=left,
  numbersep=5pt,
  numberstyle=\ttfamily\tiny\color[gray]{0.3},
  belowcaptionskip=\bigskipamount,
  captionpos=b,
  escapeinside={*'}{'*},
  language=haskell,
  tabsize=2,
  emphstyle={\bf},
  commentstyle=\it,
  stringstyle=\mdseries\rmfamily,
  showspaces=false,
  keywordstyle=\bfseries\rmfamily,
  columns=flexible,
  basicstyle=\small\sffamily,
  showstringspaces=false,
  morecomment=[l]\%,
}
% Feel free to include additional packages (if needed), but make sure it does not override the margin/font requirements.
%
\setlength{\parindent}{0pt} % no indent for new paragraphs.
\setlength{\parskip}{5pt} % distance between paragraphs, which will be used when you have a blank line in your LaTeX code.


%%%%%%%%%%%%%%%%%%%%%%%%%%%%%%%%%%%%%%%%%%%%%%%%%%%%%%%%%%%%%%%%%%%%%%%%%%%%%%
\title{Categories, Proofs, and Processes: \\Assignment 4}
%
\author{Wira Azmoon Ahmad}
\date{18 Feb 2024}
%
\begin{document}
%
\maketitle

%You may use the usual \LaTeX\ environments to structure your answers:
%\begin{align}
%        \label{basegnn}
%        \mathbf{h}_u^{(t+1)} = \sigma \bigg( \mathbf{W}_\text{self}^{(t)} \mathbf{h}_u^{(t)}  + \mathbf{W}_\text{neigh}^{(t)} \sum_{v \in N(u)} \mathbf{h}_v^{(t)} + \mathbf{b}^{(t)} \bigg): 
%\end{align}


%\begin{table}[h]
%\caption{A simple table.}
%\centering
%\begin{tabular}[t]{lrrrr}
%\toprule
%Graphs & Can & Be  & Fun & Too \\ 
%\midrule 
%$G_1$ & $1$  & $2$ & $3$ & $4$ \\ 
%$G_2$ & $-1$  & $-2$ & $-3$ & $-4$ \\ 
%$G_3$ & $0.1$  & $0.2$ & $0.3$ & $0.4$ \\ 
%\bottomrule
%\end{tabular}
%\label{tab}
%\end{table}

% A functor $F:\mathcal{C} \rightarrow \mathcal{D}$ is
% \begin{itemize}
%     \item \textbf{faithful} if each map $F_{A,B}:\mathcal{C}(A,B) \rightarrow \mathcal{D}(F\ A, F\ B)$ is injective;
%     \item \textbf{full} if each map $F_{A,B}:\mathcal{C}(A,B) \rightarrow \mathcal{D}(F\ A, F\ B)$ is surjective
% \end{itemize}

\section*{Question 1}
\addtocounter{section}{1}
\begin{proposition}
    The covariant Hom-functor $\mathcal{C}(A,-):\mathcal{C}\rightarrow$ \textbf{Set} preserves all finite limits.
\end{proposition}

\begin{proof}
Let $D: \mathcal{J} \rightarrow \mathcal{C}$ be a diagram where $\mathcal{J}$ is a finite category.
Suppose the cone $(C, (c_1,c_2,...,c_N))$ to $D$ is a finite limit.
By being a cone, for any $f:I\rightarrow J$ in $\mathcal{J}$, we have
\[D\ f\circ c_I = c_J\]

We show that $(\mathcal{C}(A,C), (\mathcal{C}(A,c_1), ,...,\mathcal{C}(A,c_N)))$ is a cone 
and a limit to $\mathcal{C}(A,-) \circ D: \mathcal{J}\rightarrow \textbf{Set}$.

Suppose we have an arbitrary $f:I\rightarrow J$ in $\mathcal{J}$,
so we get
\[(\mathcal{C}(A,-) \circ D)\ f: (\mathcal{C}(A,-) \circ D)\ I \rightarrow(\mathcal{C}(A,-) \circ D)\ J\]
which simplifies to
\[\mathcal{C}(A,(D\ f)): \mathcal{C}(A,(D\ I)) \rightarrow\mathcal{C}(A,(D\ J))\]
Consider the morphisms $D\ I \xleftarrow[]{c_I}C\xrightarrow[]{c_J} D\ J$, so
\[\mathcal{C}(A,c_I):\mathcal{C}(A,C)\rightarrow \mathcal{C}(A, (D\ I)):= h\mapsto c_I \circ h\]
\[\mathcal{C}(A,c_J):\mathcal{C}(A,C)\rightarrow \mathcal{C}(A, (D\ J)):= h'\mapsto c_J \circ h'\]
Then the diagram commutes since
\begin{equation}
\begin{aligned}
    \mathcal{C}(A,(D\ f)) \circ \mathcal{C}(A,c_I) &= (h \mapsto (D\ f)\circ h) \circ (h' \mapsto c_I \circ h')\\
    &= h \mapsto (D\ f)\circ c_I \circ h\\
    &= h \mapsto c_J \circ h\\
    &= \mathcal{C}(A,c_J)
\end{aligned}
\end{equation}
So $(\mathcal{C}(A,C), (\mathcal{C}(A,c_1), ,...,\mathcal{C}(A,c_N)))$ is a cone to $\mathcal{C}(A,-) \circ D$.
    
To show the limit, consider $\mathcal{C}(A,C)$ in \textbf{Set} and let $(B,\gamma)$ 
be some other cone to  $\mathcal{C}(A,-) \circ D$, where $\gamma$ is
\[\{\gamma_I:B\rightarrow \mathcal{C}(A,(D\ I))\}_{I\in\text{Ob}(\mathcal{J})}\]
Since $(C, (c_1,c_2,...,c_N))$ is a limit to $D$, for any cone $(A, \delta)$ to $D$, 
there is a unique arrow $a:A\rightarrow C$ such that $c_I \circ g = \delta_I$.

Note that $\forall b\in B,\gamma_I (b) \in \mathcal{C}(A,(D\ I))$.
Then for each $b\in B$, we have the cone 
\[(A,\{\gamma_I(b):A\rightarrow D\ I\}_{I\in\text{Ob}(\mathcal{J})})\]
From being the limit, there is a unique map $u_b:A\rightarrow C$ such that
\[c_I\circ u_b = \gamma_I(b)\]

By defining the arrow $g:B\rightarrow \mathcal{C}(A,C)$ such that $\forall b\in B, g(b) = u_b$,
there is a unique arrow from the cone $(B, \gamma)$ to 
$(\mathcal{C}(A,C), (\mathcal{C}(A,c_1), ,...,\mathcal{C}(A,c_N)))$, 
determined by the uniqueness of each $u_b$.

Thus, $(\mathcal{C}(A,C), (\mathcal{C}(A,c_1), ,...,\mathcal{C}(A,c_N)))$ 
is the finite limit to $\mathcal{C}(A,-) \circ D$.

    
\end{proof}

\section*{Question 2}
\addtocounter{section}{1}

\section*{Question 3}
\addtocounter{section}{1}
\section*{Question 4}
\addtocounter{section}{1}
\section*{Optional}
\addtocounter{section}{1}


\end{document}
